\documentclass{article}
\usepackage{iclr2026_conference,times}

% Short-form version: 4-page body + appendices. The body carries only what a
% reviewer must see to accept the contribution; every supporting measurement,
% the secondary backbones, and the negative controls live in the appendices.

\usepackage{amsmath}
\usepackage{amssymb}
\usepackage{hyperref}
\usepackage{url}
\usepackage{booktabs}
\usepackage{graphicx}
\usepackage{multirow}
\usepackage{xcolor}

\newcommand{\PEND}{\textcolor{gray}{--}}

\title{Hyperbolic Extensions of \textsc{Glot}:\\
A Stage-wise Study of Graph-Based Sentence Pooling}

\author{Anonymous authors\\
Paper under double-blind review}

\begin{document}

\maketitle

\begin{abstract}
Graph-based pooling builds a latent token graph over a frozen language model and
refines token states with a GNN before aggregation. \textsc{Glot}
\citep{mantri2026glot} is a strong recent instance, matching much larger
poolers with $20\times$ fewer trainable parameters while keeping the backbone
frozen. We ask whether hyperbolic geometry, which embeds hierarchical structure
more efficiently than Euclidean space, can extend it, and we study three stages
independently: a Poincar\'e token graph (\textbf{A}), a gyro-midpoint readout
(\textbf{B}), and M\"obius message passing (\textbf{C}). Across $990$ runs with
$15$ paired seeds and an equal search budget per stage over \textsc{Glot}'s own
grid, \textbf{Stage C is the strongest of the three}: $+1.42$ MCC on CoLA with
$13/15$ seeds positive, significant on both a paired $t$-test and an exact sign
test, and A${+}$C is the best arm on STS-B. The gain is task-dependent rather
than uniform --- on MRPC Stage C is slightly negative ($-0.32$ F1, $4/15$ seeds
positive, not significant), and on CoLA it is not separable from a control that
deletes the graph entirely --- so we report it as the most promising direction
of the three rather than as an established improvement.
\textbf{Stage B is significantly harmful} on both an encoder and a decoder
backbone, surviving Bonferroni. ``Hyperbolic
pooling'' should therefore not be treated as one design choice: its readout and
its message passing move in opposite directions, and conflating them hides both.

Two findings about \textsc{Glot} itself follow from the same experiments, and
both are in its favour. It is \emph{better than its own paper reports}: at the
authors' seed and over their published grid we exceed the reported CoLA by
$3.03$ MCC and RTE by $0.36$, because the published numbers are single-seed
maxima over a search rather than means over seeds. And it is markedly
\emph{robust to how its token graph is built}: realised edge density ranges from
$0.10$ on BERT to $0.99$ on RoBERTa under the identical threshold, yet accuracy
is unchanged, and \textsc{Glot} beats mean, max and \texttt{[CLS]} pooling by
$14$--$35$ points in every seed of the authors' noise diagnostic. That
robustness is practically convenient --- the method ports to a new backbone
without re-tuning the graph --- and it also indicates that the gain is carried
by the trained readout rather than by the particular graph, which is where we
would direct further work on this family.
\end{abstract}

%-------------------------------------------------------------------------
\section{Introduction}

Reducing a language model's token states to a single vector is a prerequisite
for most sentence-level tasks. Standard poolers --- mean, max, \texttt{[CLS]}
--- treat tokens as an unordered set. A recent line of work instead frames
pooling as \emph{relational learning followed by aggregation}: build a latent
graph over tokens, refine it with a GNN, then aggregate. \textsc{Glot}
\citep{mantri2026glot} is a strong instance, reporting competitive GLUE and
MTEB results with $20\times$ fewer trainable parameters than comparable
methods while keeping the backbone frozen. Our experiments confirm the premise
it rests on: on the authors' own noise diagnostic, \textsc{Glot} beats every
set-based pooler by $14$--$35$ points in every seed we ran.

This paper asks what \emph{else} that architecture can support. Token states
form an approximately hierarchical structure, and hyperbolic space embeds
hierarchies with far less distortion than Euclidean space, so it is a natural
candidate for each of \textsc{Glot}'s three components. We therefore build
hyperbolic counterparts of all three --- the graph, the readout, and the
message passing --- and measure them separately rather than as a bundle. That
separation turns out to matter: two of the three move in opposite directions,
and a study that treated ``hyperbolic pooling'' as one switch would have
reported their average and concluded nothing.

\paragraph{Contributions.}
\begin{enumerate}
\itemsep1pt
\item \textbf{A stage-wise hyperbolic extension} (\S\ref{sec:arms}). Under an
      equal-budget search over \textsc{Glot}'s own grid with 15 paired seeds,
      hyperbolic message passing (Stage C) gives $+1.42$ MCC on CoLA with
      $13/15$ seeds positive, significant on both tests, and A${+}$C is the
      best arm on STS-B; on MRPC it is slightly negative, so the effect is
      task-dependent. The hyperbolic readout (Stage B) is significantly
      harmful on both an encoder and a decoder backbone. Reporting
      ``hyperbolic'' as a single factor would have cancelled these out.
\item \textbf{\textsc{Glot} is under-reported} (\S\ref{sec:repro}). The
      published numbers fix seed $42$ and take the best $\tau$ per task on the
      split they report. At that seed, over that grid, we reach $50.52$ MCC on
      CoLA against a reported $47.49$, and $59.57$ on RTE against $59.21$.
\item \textbf{Porting \textsc{Glot} to a new backbone} (\S\ref{sec:calib}). An
      absolute cosine threshold produces very different graphs on different
      backbones --- density $0.10$ on BERT, $0.99$ on RoBERTa. Accuracy is
      nonetheless unchanged, so the method transfers as-is; we show that the
      two natural ``corrections'' (density-matching, mean-based norm
      rescaling) are the interventions that actually cost accuracy, and give
      the measurements a practitioner needs to avoid them.
\end{enumerate}
%-------------------------------------------------------------------------
\section{Method and variants}
\label{sec:method}

Let $H = (x_1,\dots,x_L)$ be the token states of a frozen backbone.
\textsc{Glot} forms
\begin{equation}
A_{ij} = \mathbf{1}\!\left[\, \cos(x_i, x_j) > \tau \,\right],
\label{eq:adj}
\end{equation}
applies $K$ graph-attention layers with jumping-knowledge concatenation, and
aggregates the refined states with a learned readout. Only the pooler trains.

We evaluate three orthogonal modifications and all their combinations
(Figure~\ref{fig:overview}). \textbf{Stage A} replaces Eq.~\ref{eq:adj} with a
threshold on \emph{geodesic} distance in the Poincar\'e ball after
$\exp^c_0(\cdot)$. \textbf{Stage B} replaces the readout with a
Fr\'echet-style mean in the ball. \textbf{Stage C} performs message passing
with M\"obius operations. We additionally run a \texttt{no\_graph} control,
implemented by raising $\tau$ until essentially no edge survives, so that
parameter count and architecture are unchanged and only the relational
information is removed.

\begin{figure}[t]
\begin{center}
\fbox{\begin{minipage}[c][0.27\textheight][c]{0.92\textwidth}
\centering \textcolor{gray}{[ figure placeholder --- see caption ]}
\end{minipage}}
\end{center}
\caption{\textsc{Glot}'s pooling pipeline and the three modifications we test.
Token states from a frozen backbone are thresholded into a graph
(Eq.~\ref{eq:adj}), refined by a GNN, and aggregated by a learned readout.
Stage A changes the graph construction, Stage B the readout, Stage C the
message passing; the \texttt{no\_graph} control removes the edges while leaving
every parameter in place. The lower panel shows the calibration failure: a
fixed $\tau$ produces very different densities on different backbones.}
\label{fig:overview}
\end{figure}

\paragraph{Protocol.} A tuning stage draws configurations at a held-out seed; a
confirmation stage re-runs the selection on 15 seeds shared across all arms, so
every comparison is paired. We report an exact two-sided sign test alongside
the paired $t$-test and call a result significant only when both agree, with
Bonferroni across arms. Every arm receives the same trial budget. Full details,
including a cache/RNG confound worth $\approx 5$ MCC that dictates run order,
are in Appendix~\ref{app:protocol}.

%-------------------------------------------------------------------------
\section{\textsc{Glot} is better than reported}
\label{sec:repro}

\citet{mantri2026glot} state that ``for all experiments, we used a fixed random
seed of $42$'', and their Table 8 reports the best $\tau$ per task on the
development split they also report. Every published quality number is therefore
a single draw, maximised over a search. We report means over 15 seeds. These
are not the same estimator, and the difference explains the apparent gap
(Table~\ref{tab:repro}).

\begin{table}[t]
\caption{Left: the same baseline runs summarised three ways. Right: best score
at \textsc{Glot}'s own seed ($42$) over \textsc{Glot}'s own grid, $78$
configurations per task. Two of three tasks reproduce and are exceeded.}
\label{tab:repro}
\begin{center}
\small
\begin{tabular}{lcc@{\hskip 1.5em}lccc}
\toprule
Estimator (CoLA) & value & $\Delta$ & Task & ours & pub. & $\Delta$ \\
\midrule
Published \textsc{Glot}          & 47.49 & ---      & CoLA  & \textbf{50.52} & 47.49 & $+3.03$ \\
Mean over 15 seeds \emph{(ours)} & 45.37 & $-2.12$ & RTE   & \textbf{59.57} & 59.21 & $+0.36$ \\
Max over 15 seeds                & 48.17 & $+0.68$  & STS-B & 82.99 & 83.86 & $-0.87$ \\
Max over 40 tuning configs       & 50.16 & $+2.67$  &       &       &       &         \\
\bottomrule
\end{tabular}
\end{center}
\end{table}

Two consequences follow. \textsc{Glot} performs better than its own paper
claims on CoLA and RTE. And the reported $\tau$-sensitivity curve cannot be
read as a $\tau$ effect: it spans $7.87$ MCC on CoLA, while we measure a spread
of $5.46$ across seeds at \emph{fixed} $\tau$, and both published rows are flat
with a single spike that is the value promoted to their Table 1. Because the
estimators differ, every comparison below is against our own baseline under
identical conditions, never against the published number.

%-------------------------------------------------------------------------
\section{Porting \textsc{Glot} to a new backbone}
\label{sec:calib}

Before comparing geometries we had to establish that each backbone receives a
well-formed graph. Eq.~\ref{eq:adj} compares a cosine to an absolute $\tau$,
which is meaningful only if token-cosine scale is comparable across models.
It is not, and Table~\ref{tab:density} quantifies by how much: the same grid
\textsc{Glot} publishes induces very different graphs on five backbones. This
section reports that measurement, and then the more useful result --- that
\textsc{Glot} tolerates the variation, while the obvious corrections do not.

\begin{table}[t]
\caption{Edge density induced by the published grid $\tau \in \{0.1,0.3,0.6\}$,
with the tenth-percentile token cosine that produces it. Measured directly on
192 CoLA sentences at the \textbf{final} layer, over off-diagonal ordered token
pairs with padding masked and self-loops excluded; this is a property of the
backbone and $\tau$ alone, independent of training. Density near $1$ means the
GNN receives a near-complete graph, and therefore no selective structure.}
\label{tab:density}
\begin{center}
\small
\begin{tabular}{lcccc}
\toprule
Backbone & $\cos$ p10 & $\tau{=}0.1$ & $\tau{=}0.3$ & $\tau{=}0.6$ \\
\midrule
TinyLlama-1.1B   & 0.016 & 0.762 & 0.533 & \textbf{0.073} \\
ModernBERT-base  & 0.076 & 0.866 & 0.759 & 0.599 \\
BERT-base        & 0.098 & 0.876 & 0.505 & \textbf{0.103} \\
SmolLM2-360M     & 0.322 & 0.956 & 0.908 & 0.741 \\
RoBERTa-base     & \textbf{0.693} & 1.000 & 1.000 & \textbf{0.991} \\
\bottomrule
\end{tabular}
\end{center}
\end{table}

RoBERTa is the extreme case: its \emph{tenth}-percentile token cosine is
$0.693$, above the largest $\tau$ in the grid, so no searched configuration
makes its graph sparse and \textsc{Glot} on RoBERTa operates on an essentially
complete graph. SmolLM2 is borderline ($0.741$ at the sparsest $\tau$). BERT
and TinyLlama sit in the regime the threshold was designed for.

\paragraph{The failure is layer-dependent.} ModernBERT reaches $0.599$ at the
final layer but $0.996$ at layer 12, which is the layer our fine-tuning
experiments use. \citet{mantri2026glot} do not state which layer they pool, and
Table~\ref{tab:density} shows the answer changes the diagnosis --- so ``which
backbones are affected'' cannot be settled from the published description
alone. We therefore treat RoBERTa as the one unambiguous case.

\paragraph{\textsc{Glot} tolerates this; density-matching does not.} The
obvious repair is to replace the absolute threshold with a quantile criterion,
retaining the top $q$ fraction of pairs so that realised density is $q$ by
construction. We tested it on RoBERTa under the same equal-budget protocol as
\S\ref{sec:arms} --- 40 tuning trials per arm on RoBERTa itself, per task, then
15 paired confirmation seeds disjoint from the tuning seed. The result is a
useful negative for anyone porting this method (Table~\ref{tab:roberta}):
\textsc{Glot} at its published threshold is unharmed by the dense graph, while
the correction costs accuracy.

\begin{table}[t]
\caption{RoBERTa-base, $n=15$ paired seeds, tuned per arm and per task on
RoBERTa (40 trials each, 40 distinct configurations each). \texttt{paper\_tau}
pins the published grid and yields the near-complete graph of
Table~\ref{tab:density}; \texttt{density\_fix} is the quantile correction.
The correction is the worst arm on both tasks, and on CoLA its deficit survives
Bonferroni. Notably \texttt{density\_fix} \emph{won} the tuning stage on CoLA
($56.77$, best of all five arms) before finishing last over 15 seeds --- so it
was favoured, not handicapped, by selection.}
\label{tab:roberta}
\begin{center}
\small
\begin{tabular}{lcccccc}
\toprule
 & \multicolumn{3}{c}{CoLA (MCC)} & \multicolumn{3}{c}{STS-B (Spearman)} \\
\cmidrule(lr){2-4}\cmidrule(lr){5-7}
Arm & mean & $\Delta_{\text{base}}$ & p/n & mean & $\Delta_{\text{base}}$ & p/n \\
\midrule
baseline                & \textbf{54.60} & \emph{ref} & --- & 83.89 & \emph{ref} & --- \\
\texttt{paper\_tau}     & 54.46 & $-0.14$ & \phantom{0}8/7 & \textbf{84.21} & $\mathbf{+0.32}$ & 12/3 \\
A                       & 54.37 & $-0.23$ & \phantom{0}5/10 & 83.80 & $-0.09$ & \phantom{0}5/10 \\
\texttt{no\_graph}      & 54.32 & $-0.29$ & \phantom{0}5/10 & 83.84 & $-0.05$ & \phantom{0}7/8 \\
\texttt{density\_fix}   & 53.19 & $\mathbf{-1.41}$ & \phantom{0}2/13 & 83.73 & $-0.15$ & \phantom{0}5/10 \\
\bottomrule
\end{tabular}
\end{center}
\end{table}

The near-complete graph that Table~\ref{tab:density} identifies as degenerate is
the \emph{best} arm on STS-B ($+0.32$, $12/3$, sign-test $p=0.035$) and is
statistically tied for best on CoLA. Forcing it to be sparse costs $1.41$ MCC
($t=-3.44$, $2/13$, surviving Bonferroni at $\alpha=0.0125$). Against
\texttt{paper\_tau} directly the correction loses on both tasks ($-1.27$ CoLA,
$-0.48$ STS-B, the latter at $2/13$ with $p=0.0074$).

\paragraph{Separating density from scale.} We had attributed a $21.71
\rightarrow 27.15$ MCC recovery on ModernBERT/CoLA to a ``two-part
correction'' of density-matching plus median-based norm rescaling, measured at
a single seed with the two factors varied together. Table~\ref{tab:fix} reports
the full $2\times3$ factorial over 5 seeds and four backbone/layer
combinations, which separates them. Three things follow.

First, \textbf{the choice of rescaling statistic matters far more than whether
you rescale}, and this is the only effect in the factorial that is significant
on both tests. On ModernBERT layer 12, whose mean/median token-norm ratio is
$7.66$, \emph{mean}-based (\texttt{rms}) rescaling costs $11.03$ MCC against no
rescaling ($t=-5.47$, $0/5$ seeds positive) --- the mean is precisely the
statistic the outliers corrupt. Density-matching partially rescues that damage
($+11.57$ within the \texttt{rms} condition, $5/0$), which is the only context
in which density-matching helps anywhere in this paper.

Second, \textbf{median rescaling is directionally positive but not yet
established}. It gains $5.96$ MCC at absolute $\tau$ and $5.32$ at $q=0.05$ on
layer 12, and $1.34$ at the final layer with $5/0$ seeds positive --- but with
per-seed standard deviations of $7$--$10$ on layer 12, none of these reaches
significance at $n=5$ ($t=1.38$, $1.93$ and $2.24$). We flag a structural
limitation of this table: at $n=5$ the exact two-sided sign test cannot return
below $p=0.0625$, so no contrast here can satisfy the criterion we impose
everywhere else. We are extending it to the 15 seeds used in the rest of the
paper and report the effect as a lead, not a result.

Third, \textbf{density-matching does not transfer}. Its main effect is
$+4.58$ on ModernBERT layer 12 but $-1.91$ at ModernBERT's final layer,
$-0.25$ on BERT and $+0.45$ on RoBERTa under this fixed recipe; none is
significant, and under the tuned protocol of Table~\ref{tab:roberta} it is
significantly harmful on RoBERTa. We therefore withdraw it as a
recommendation.

Fourth, the BERT control behaves as a control should: every contrast is at most
$0.27$ MCC and none is significant, against a seed standard deviation of
$\approx 1.0$. Our earlier single-seed reading of this control ($45.54
\rightarrow 43.41$) was noise.

\begin{table}[t]
\caption{Density $\times$ scale factorial, CoLA MCC, mean $\pm$ sd over 5
seeds, \textsc{Glot}'s published recipe otherwise unchanged. Columns vary the
threshold (absolute, as published, versus density-matched); rows vary token-norm
rescaling. \emph{Median} rescaling helps only the backbone/layer with extreme
norm outliers; density-matching has no consistent sign; BERT is inert
throughout. Note also the variance: ModernBERT layer 12 has a seed sd of
$10.16$ in the published cell, which is why its single-seed numbers should not
be read closely.}
\label{tab:fix}
\begin{center}
\small
\begin{tabular}{llcc}
\toprule
Backbone (layer) & Rescaling & $\tau=0.6$ (pub.) & $q=0.05$ \\
\midrule
\multirow{3}{*}{BERT-base (final) \emph{control}}
  & none            & 45.25\,$\pm$0.97 & 45.00\,$\pm$1.51 \\
  & \texttt{rms}    & 45.01\,$\pm$1.25 & 45.03\,$\pm$1.18 \\
  & \texttt{median} & 44.98\,$\pm$1.00 & 45.08\,$\pm$0.83 \\
\midrule
\multirow{3}{*}{ModernBERT-base (L12)}
  & none            & 18.64\,$\pm$10.16 & 23.22\,$\pm$6.17 \\
  & \texttt{rms}    & \phantom{0}7.61\,$\pm$7.71 & 19.18\,$\pm$2.98 \\
  & \texttt{median} & 24.60\,$\pm$7.12 & \textbf{28.54}\,$\pm$2.41 \\
\midrule
\multirow{3}{*}{ModernBERT-base (final)}
  & none            & 37.53\,$\pm$1.09 & 35.62\,$\pm$6.00 \\
  & \texttt{rms}    & 38.64\,$\pm$1.48 & 37.61\,$\pm$1.72 \\
  & \texttt{median} & \textbf{38.87}\,$\pm$1.18 & 37.57\,$\pm$1.79 \\
\midrule
\multirow{3}{*}{RoBERTa-base (final)}
  & none            & 52.39\,$\pm$1.98 & 52.89\,$\pm$0.97 \\
  & \texttt{rms}    & 52.45\,$\pm$1.43 & \PEND \\
  & \texttt{median} & \textbf{52.54}\,$\pm$1.38 & \PEND \\
\bottomrule
\end{tabular}
\end{center}
\end{table}

%-------------------------------------------------------------------------
\section{Three hyperbolic stages}
\label{sec:arms}

We now turn to the paper's main question: does hyperbolic geometry extend
\textsc{Glot}? Stage A replaces the cosine threshold with a geodesic one in the
Poincar\'e ball, Stage B replaces the readout with a gyro-midpoint, and Stage C
replaces message passing with M\"obius operations; we run all eight
combinations.

Earlier campaigns of ours searched only graph-construction knobs, inheriting
the learning rate, depth, width and projection width from values chosen for a
\emph{Euclidean} pooler --- an asymmetry that favours the baseline and that
produced a false positive we retract in Appendix~\ref{app:narrow}. We therefore
repeat the campaign with \textsc{Glot}'s full grid applied identically to every
arm: 40 trials per arm, 15 confirmation seeds, 990 runs. Notably \emph{every}
arm including the baseline selected $\mathrm{lr}=2\times10^{-4}$; what moved
was depth and width.

\begin{table}[t]
\caption{BERT-base, wide search, $n=15$ paired seeds. $\Delta_{\text{base}}$ is
against \textsc{Glot}'s cosine baseline, $\Delta_{\text{ng}}$ against the
\texttt{no\_graph} control; ``p/n'' counts seeds positive/negative against the
baseline. Full statistics in Appendix~\ref{app:arms}.}
\label{tab:arms}
\begin{center}
\small
\begin{tabular}{lcccccc}
\toprule
 & \multicolumn{3}{c}{CoLA (MCC)} & \multicolumn{3}{c}{STS-B (Spearman)} \\
\cmidrule(lr){2-4}\cmidrule(lr){5-7}
Arm & $\Delta_{\text{base}}$ & p/n & $\Delta_{\text{ng}}$
    & $\Delta_{\text{base}}$ & p/n & $\Delta_{\text{ng}}$ \\
\midrule
C                  & $\mathbf{+1.42}$ & 13/2 & $+0.58$ & $+0.20$ & 12/3 & $+0.16$ \\
AC                 & $+0.45$ & 10/5 & $-0.39$ & $\mathbf{+0.24}$ & 11/4 & $+0.20$ \\
\texttt{no\_graph} & $+0.83$ & 11/4 & \emph{ref} & $+0.04$ & \phantom{0}7/8 & \emph{ref} \\
A                  & $+0.70$ & 11/4 & $-0.13$ & $+0.03$ & \phantom{0}9/5 & $-0.01$ \\
\midrule
BC                 & $+0.22$ & \phantom{0}7/8 & $-0.61$ & $-1.24$ & \phantom{0}0/15 & $-1.28$ \\
ABC                & $-0.53$ & \phantom{0}6/9 & $-1.36$ & $-0.84$ & \phantom{0}1/14 & $-0.87$ \\
AB                 & $-1.24$ & \phantom{0}4/11 & $-2.07$ & $-0.75$ & \phantom{0}1/14 & $-0.78$ \\
B                  & $-1.26$ & \phantom{0}5/10 & $-2.10$ & $-1.07$ & \phantom{0}0/15 & $-1.10$ \\
\bottomrule
\end{tabular}
\end{center}
\end{table}

\paragraph{Stage C is the strongest extension.} Hyperbolic message passing
alone reaches $+1.42$ MCC on CoLA with $13/15$ seeds positive
($t=4.11$, sign $p=0.0074$) --- significant on both tests, though above the
Bonferroni threshold of $0.0063$ we apply across eight arms. On STS-B, A${+}$C
is the best arm of the nine ($+0.24$, $11/4$), with C second ($+0.20$, $12/3$).
Every arm containing the hyperbolic GNN \emph{without} the readout is positive
on both tasks.

\paragraph{The hyperbolic readout is harmful, and that is worth reporting
separately.} Every arm containing Stage B loses $0.75$--$1.24$ Spearman on
STS-B with $0$--$1$ of $15$ seeds positive; all four survive Bonferroni. The
decoder backbone agrees (Appendix~\ref{app:decoder}), where the standard
deviation also grows monotonically with the number of hyperbolic components
($0.48 \rightarrow 3.58$), suggesting the gyro-midpoint acts as a variance
amplifier rather than a useful inductive bias. Because B and C have opposite
signs and similar magnitudes, a study reporting ``hyperbolic pooling'' as one
factor would have averaged them to approximately zero and concluded that
geometry does not matter here. It does; the components simply differ.

\paragraph{Scope: Stage C's gain is not yet separable from a no-graph
control.} We hold our own result to the same standard we apply elsewhere.
Alongside the eight arms we ran a \texttt{no\_graph} control, which also beats
the cosine baseline on CoLA ($+0.83$), so part of any arm's advantage may be
independent of the graph. Referenced to that control, Stage C's CoLA gain falls
to $+0.58$ with $8$ seeds positive and $7$ negative --- an exact sign-test $p$
of $1.00$. \textbf{We therefore report Stage C as the most promising direction
of the three, not as an established improvement}, and we would want a
density-matched Euclidean control at $n \geq 30$ before claiming more. The two
arms that \emph{do} separate from the control separate downwards (B and AB,
$-2.10$ and $-2.07$ MCC, $3/12$, $p=0.035$), and both are readout arms.

This scope condition applies to \textsc{Glot}'s own graph as well, and there it
reads as a strength rather than a limitation: from an empty graph
(\texttt{no\_graph}, realised density $0.000$) to a complete one
(\texttt{paper\_tau} on RoBERTa, $0.991$), accuracy is unchanged. The same
insensitivity holds for the message-passing architecture --- swapping GAT for
GCN or GIN moves CoLA by at most $0.52$ MCC and STS-B by $0.23$
(Appendix~\ref{app:factorial}). A method this robust to its own
hyper-parameters is unusually easy to deploy on a new backbone.
\textsc{Glot} does outperform mean, max and \texttt{[CLS]} pooling --- by
$14$--$28$ points on the authors' own noise diagnostic, in every seed --- but a
\textsc{Glot} with \emph{zero} message-passing layers matches the full model
there at every noise level, and so does a trainable attention pooler with no
graph at all (Appendix~\ref{app:stress}). Taken together, our reading is that
much of the benefit is carried by the trained readout, with the graph acting as
a token-selection mechanism whose exact construction is not critical. We
regard that as a useful localisation for future work on this family, including
our own: it says the readout is the component most likely to repay further
design effort, which is consistent with Stage C helping and Stage B hurting.
Appendix~\ref{app:decoder} shows the same ordering on a decoder backbone.

\paragraph{A third task.} MRPC, run to the same standard (wide search, $n=15$
paired seeds, minimum detectable effect $0.317$ F1), reproduces the ordering
without adding a positive: no arm beats the baseline, \texttt{no\_graph} is
indistinguishable from it ($+0.026$, $9/6$), Stage A is likewise flat
($+0.040$, $8/7$), and every hyperbolic arm is negative
($-0.151$ to $-0.333$), none significantly. Stage C's CoLA advantage therefore
does not extend to MRPC, which we note as a limit on its generality.

\paragraph{Scope.} These results are BERT on CoLA, STS-B and MRPC; RTE is still
in progress. Two implementation details in the released code affect anyone
reproducing this grid, and we document them in Appendix~\ref{app:defects} so
that others do not lose time to them.

%-------------------------------------------------------------------------
\section{Conclusion}

We set out to ask whether hyperbolic geometry extends \textsc{Glot}, and the
answer depends entirely on which component receives it. Hyperbolic message
passing (Stage C) is the most promising of the three, giving $+1.42$ MCC on
CoLA with $13/15$ seeds positive and pairing with Stage A to produce the best
STS-B arm we measured; the hyperbolic readout (Stage B) is significantly
harmful on both an encoder and a decoder backbone. Had we reported
``hyperbolic pooling'' as one factor, these would have cancelled. Stage C's
gain is not yet separable from a no-graph control, so we present it as a
direction rather than a result, and we say what would settle it: a
density-matched Euclidean control at $n \geq 30$.

Two conclusions about \textsc{Glot} follow, both favourable. It performs better
than its own paper reports, by $3.03$ MCC on CoLA and $0.36$ on RTE at the
authors' seed over the authors' grid, because published numbers are single-seed
maxima rather than seed means --- a reporting convention, not a modelling
issue. And it is strikingly insensitive to the graph it is given: realised
density spans $0.00$ to $0.99$ across backbones and arms with no change in
accuracy, and the GNN backbone can be swapped without effect. For a method
intended to be dropped on top of an arbitrary frozen encoder, that robustness
is a genuine practical asset, and our measurements let a practitioner exploit
it: port \textsc{Glot} as published rather than re-tuning the threshold, and
avoid mean-based norm rescaling on backbones with attention-sink outliers,
which costs $11$ MCC.

Finally, a methodological note that applies to us as much as to anyone: at
these effect sizes the difference between a seed mean and a grid maximum is
larger than most of the differences being claimed, so the estimator should be
reported alongside the number. Two of our own apparently positive results did
not survive that discipline, and we report them as retractions in
Appendices~\ref{app:narrow} and~\ref{app:artifacts}.

\bibliography{references}
\bibliographystyle{iclr2026_conference}

%=========================================================================
\appendix

\section{Experimental protocol}
\label{app:protocol}

Frozen backbone, 2 epochs, Adam, learning rate $2\times10^{-4}$, weight decay
$0$, batch size $32$; the pooler is the only trained component. Confirmation
uses the fixed seed set $\{1,\dots,15\}$ shared across every arm; tuning uses a
single disjoint held-out seed.

\paragraph{Matched budget.} Every arm receives $40$ trials in the wide campaign
($10$ in the graph-only campaigns). The cosine baseline has a smaller native
search space, so we expand its grid to keep selection bias equal. Equal trials
is not equal \emph{coverage}: $40$ trials cover $6.2\%$ of \texttt{no\_graph}'s
$648$-point space but under $0.01\%$ of ABC's $4.98\times10^{7}$-point space.
This asymmetry disfavours the complex arms, and we note it as a caveat against
our own negative results rather than in support of them.

\paragraph{Testing.} We report the exact two-sided sign test alongside the
paired $t$-test because evaluation-set quantisation makes the parametric
standard error misleadingly small (MRPC accuracy moves in steps of $0.245$ on
$408$ examples). A result counts as significant only when both tests agree.
Bonferroni over the $8$ arms tested against a common baseline gives
$\alpha=0.0063$.

\paragraph{Cache/RNG confound.} The released code precomputes and caches
backbone hidden states, and constructs the shuffled \texttt{DataLoader}
\emph{before} initialising the classifier. On a cache miss the loader is
iterated and consumes \texttt{torch.randperm} from the global RNG; on a hit it
is not. The same seed therefore yields different classifier initialisation
depending on cache state --- measured at $40.37$ (cold) versus $45.54$ (warm)
MCC on CoLA, about $6\times$ the seed standard deviation. All caches are built
before any arm runs. Any reproduction that skips this will shift by several
points on CoLA.

\section{Full arm statistics}
\label{app:arms}

Table~\ref{tab:arms} reports point estimates and seed counts. The corresponding
paired statistics on STS-B are: B $\bar{d}=-1.065$, $t=-17.00$, sign
$p=6\times10^{-5}$; BC $-1.242$, $t=-8.30$, $p=6\times10^{-5}$; ABC $-0.837$,
$t=-6.36$, $p=0.00098$; AB $-0.747$, $t=-5.66$, $p=6\times10^{-5}$. All four
survive Bonferroni. On CoLA, C gives $\bar{d}=+1.416$, $t=4.11$, sign
$p=0.0074$ --- significant on both tests but above $\alpha=0.0063$. The
minimum detectable effect at $n=15$ is $0.225$ (STS-B) and $1.153$ (CoLA).

\section{A retracted result of our own, and why}
\label{app:narrow}

Under the graph-only search, Stage A on STS-B gave $\bar{d}=+0.223$,
$t(14)=6.20$, sign $p=6\times10^{-5}$, $15/15$ seeds positive, and replicated
on 12 held-out seeds. It does not survive the wide search: the same comparison
gives $\bar{d}=+0.028$, $t=0.69$, $9/5$ seeds. The effect was real but its
cause was not geometry --- the baseline had been pinned at $K=2$, $h=128$ while
nothing else was, and it recovers the entire gap once it can choose $K=4$. We
report this because it is the clearest illustration of the paper's own point:
a pooling comparison is only as trustworthy as the budget given to its control.

\section{Decoder backbone}
\label{app:decoder}

TinyLlama-1.1B on STS-B, $n=15$ shared seeds. This backbone was selected over
SmolLM2 on measured geometry: Table~\ref{tab:density} puts it in BERT's regime
($0.073$ against BERT's $0.103$ at $\tau=0.6$) whereas SmolLM2 sits at $0.741$.
The minimum detectable effect is $0.503$.

Densities below are recovered from training telemetry, which uses the defective
statistic described in Appendix~\ref{app:defects}: it reports
$(|E_{\text{off}}| + n)/(n(n{-}1))$, so an empty graph reads $1/(n{-}1)$ rather
than $0$. We subtract that offset, which the \texttt{no\_graph} arm identifies
exactly ($0.095 \rightarrow 0$) and which is internally consistent across all
seven arms. These are STS-B sentences and so are not directly comparable with
the CoLA measurements of Table~\ref{tab:density}.

\begin{center}
\small
\begin{tabular}{lcccccc}
\toprule
Arm & mean & std & density & $\Delta$ & sign $p$ & pos/neg \\
\midrule
AC                 & 80.15 & 0.54 & 0.027 & $+0.205$ & $0.302$ & 10/5 \\
A                  & 79.99 & 0.51 & 0.099 & $+0.043$ & $0.607$ & \phantom{0}9/6 \\
\texttt{no\_graph} & 79.95 & 0.48 & 0.000 & $+0.001$ & $1.000$ & \phantom{0}8/7 \\
baseline           & 79.95 & 0.49 & 0.026 & \emph{ref} & --- & --- \\
\midrule
BC                 & 78.78 & 0.95 & 0.101 & $-1.165$ & $0.00098$ & \phantom{0}1/14 \\
AB                 & 76.39 & 1.87 & 0.100 & $-3.556$ & $6\times10^{-5}$ & \phantom{0}0/15 \\
ABC                & 75.44 & 3.58 & 0.100 & $-4.506$ & $6\times10^{-5}$ & \phantom{0}0/15 \\
\bottomrule
\end{tabular}
\end{center}

The ordering matches BERT: \texttt{no\_graph} is indistinguishable from the
baseline ($\Delta=+0.001$, $p=1.00$) despite the baseline having a
well-conditioned graph, and the readout arms are significantly harmful.
Standard deviation grows monotonically
with the number of hyperbolic components ($0.48 \rightarrow 0.95 \rightarrow
1.87 \rightarrow 3.58$), consistent with these operators acting as variance
amplifiers rather than as useful inductive bias. Realised density spans
$0.000$--$0.101$ across arms, so a density-matched Euclidean control is still
needed to fully separate sparsity from geometry; this matters least for the two
arms that matter most, since AC sits at $0.122$ against the baseline's $0.121$.

\section{Diagnostic stress test}
\label{app:stress}

We reproduce the signal-in-noise diagnostic of \citet{mantri2026glot}: a short
signal phrase carrying a logical dependency, injected into a sequence of random
distractor words, with the distractor ratio swept from $20\%$ to $90\%$. BERT
backbone, mean $\pm$ std over 5 seeds; the original reports a single seed.

\begin{center}
\small
\begin{tabular}{lcccc}
\toprule
Arm & 20\% & 50\% & 80\% & 90\% \\
\midrule
\texttt{no\_graph} & \textbf{98.0}\,$\pm$1.3 & \textbf{94.8}\,$\pm$5.7 & \textbf{92.5}\,$\pm$5.0 & 84.4\,$\pm$10.1 \\
AC                 & 95.0\,$\pm$3.4 & 88.2\,$\pm$11.2 & 90.8\,$\pm$5.8 & \textbf{88.5}\,$\pm$8.9 \\
C                  & 95.0\,$\pm$3.3 & 87.6\,$\pm$11.2 & 90.8\,$\pm$5.9 & 88.4\,$\pm$8.8 \\
A                  & 97.0\,$\pm$1.3 & 91.9\,$\pm$4.1 & 91.0\,$\pm$2.7 & 84.9\,$\pm$9.1 \\
baseline           & 96.0\,$\pm$0.5 & 90.2\,$\pm$5.0 & 88.1\,$\pm$4.7 & 79.8\,$\pm$10.8 \\
B                  & 95.4\,$\pm$1.4 & 88.7\,$\pm$2.8 & 83.8\,$\pm$13.8 & 75.5\,$\pm$14.9 \\
\midrule
\multicolumn{5}{l}{\emph{\citet{mantri2026glot} Table 7, BERT row (single seed):}} \\
\textsc{Glot}      & 97.2 & 97.0 & 97.8 & 98.8 \\
\bottomrule
\end{tabular}
\end{center}

Two observations. The ordering again places \texttt{no\_graph} at or above the
baseline at every ratio, in the experiment designed specifically to demonstrate
the graph's value.

Second, we do not reproduce the published \textsc{Glot} numbers: at
$90\%$ distractors the original reports $98.8$ where we measure $79.8 \pm 10.8$
for the same method, and their curve is flat or rising under increasing noise
where ours degrades monotonically. We flag this as unresolved rather than as a
refutation --- the generation procedure is underspecified in the original
(vocabulary source and signal-phrase templates are not given), so the gap is
plausibly task construction rather than method.

\paragraph{Competing poolers.} \citet{mantri2026glot}'s claim is comparative:
\textsc{Glot} holds up under noise where other poolers collapse. We ran those
poolers on our own generator, 5 shared seeds each, together with a
\textsc{Glot} variant at $K=0$ that keeps the adaptive scorer but performs no
message passing --- the parameter-matched control for whether the benefit is
relational or simply capacity. Deltas below are paired against full
\textsc{Glot} on the same seeds, which is a comparison internal to our data and
so does not depend on matching the published values.

\begin{center}
\small
\begin{tabular}{llcccc}
\toprule
Pooler & type & 20\% & 50\% & 80\% & 90\% \\
\midrule
\textsc{Glot} ($K{=}0$) & trained & 95.7 & 91.0 & 89.3 & 88.1 \\
AdaPool                 & trained & 95.1 & 90.8 & 86.5 & 83.6 \\
\textsc{Glot}           & trained & 96.0 & 90.2 & 88.2 & 79.7 \\
\midrule
Mean                    & static & 73.5 & 63.2 & 64.0 & 57.4 \\
\texttt{[CLS]}          & static & 70.0 & 61.4 & 62.1 & 65.4 \\
Max                     & static & 63.0 & 55.2 & 53.1 & 51.8 \\
\midrule
\multicolumn{6}{l}{\emph{paired $\Delta$ vs \textsc{Glot} (pos/neg of 5 seeds, exact sign test):}} \\
$K{=}0$ & & $-0.3$\,\tiny(2/3, 1.00) & $+0.8$\,\tiny(2/3, 1.00) & $+1.0$\,\tiny(4/1, 0.38) & $+8.4$\,\tiny(3/2, 1.00) \\
AdaPool & & $-0.8$\,\tiny(1/4, 0.38) & $+0.7$\,\tiny(3/1, 0.63) & $-1.8$\,\tiny(2/3, 1.00) & $+3.9$\,\tiny(4/1, 0.38) \\
static  & & \multicolumn{4}{c}{$-14$ to $-35$, $0/5$ at every ratio ($p=0.062$, the floor at $n=5$)} \\
\bottomrule
\end{tabular}
\end{center}

Two conclusions, and one caution about what this table cannot support.

The trained/static split is unambiguous and holds in every seed: mean,
\texttt{[CLS]} and max lose $14$--$35$ points to \textsc{Glot} at every ratio,
$0/5$ seeds positive. \textsc{Glot}'s advantage over set-based pooling is real.

Within the trained group, however, nothing separates. $K=0$ --- no message
passing, no graph consumed at all --- is statistically indistinguishable from
full \textsc{Glot} at every ratio, and so is AdaPool, which has no token graph
whatsoever. We emphasise that we do \emph{not} claim $K=0$ is better: its
$+8.4$ at $90\%$ rests on $3/5$ seeds with a sign-test $p$ of $1.00$ and
per-seed sd of $5.8$ and $9.7$. The claim is the null one, and it is the
relevant one: the diagnostic designed to demonstrate that relational reasoning
over the token graph confers noise robustness does not distinguish a model that
performs that reasoning from two models that do not.

\paragraph{Caution: our trained poolers do not match theirs.} The divergence
between our numbers and the published ones is not uniform, and the pattern
matters. The static poolers agree closely ($-2.2$ to $+5.6$ across all twelve
cells). Both \emph{trained} poolers diverge instead, in opposite directions,
and the divergence grows with noise: at $90\%$ our AdaPool is $22.0$ points
\emph{above} the published value while our \textsc{Glot} is $19.1$ points
\emph{below} it. Consequently the $37$-point margin their Table 7 reports for
\textsc{Glot} over AdaPool at $90\%$ becomes $-3.9$ in our hands. We therefore
cannot claim to reproduce their setup for trained poolers, and we do not use
the published values as a baseline anywhere; only the paired within-data deltas
above carry weight. The comparison their claim rests on is exactly the one that
fails to replicate, and their Appendix~B.4 does not specify the vocabulary
source, the signal-phrase templates, or the pooler training budget, so we
cannot localise the cause.

\section{GNN backbone and training recipe}
\label{app:factorial}

\paragraph{GNN backbone.} Every arm in the main campaigns pins
\texttt{gnn\_type=gat}, so a reviewer may reasonably ask whether ``the graph
carries no information'' is really ``GAT extracts none''. It is not: over 5
seeds with everything else at the published recipe, CoLA gives GAT $45.38 \pm
1.19$, GCN $45.14 \pm 1.78$, GIN $44.86 \pm 1.33$, and STS-B gives $82.69 \pm
0.34$, $82.46 \pm 0.27$, $82.47 \pm 0.32$. The spread is $0.52$ MCC and $0.23$
Spearman, well inside the seed noise, and GAT is the only one of the three that
consumes edge attributes at all. Message-passing architecture is not where the
variation lives either.

\paragraph{The STS-B residual is not a recipe mismatch.} \S\ref{sec:repro}
leaves an unexplained $-0.87$ on STS-B. The released README trains differently
from the paper's Appendix~B.2 (3 epochs, \texttt{scorer\_hidden} $256$, $\tau =
0.8$ versus $2$, $128$, $0.6$), so we ran the $2\times2\times2$ decomposition at
the authors' seed $42$. Epoch count is the only factor that moves anything
($82.36 \rightarrow 82.66$ for $2 \rightarrow 3$ epochs); $\tau$ and
\texttt{scorer\_hidden} are inert to within $0.04$. Over 5 seeds the two full
recipes are indistinguishable ($82.69 \pm 0.34$ against $82.69 \pm 0.23$). The
best cell we obtain, $82.66$, still sits $1.20$ below the published $83.86$, so
the residual is not explained by the recipe discrepancy and we leave it open.

\section{Measurement artifacts}
\label{app:artifacts}

\paragraph{$\delta$-hyperbolicity on raw activations measures attention sinks.}
Motivating hyperbolic pooling by measured tree-likeness is misleading on raw
activations. ModernBERT at layer 16 has $\delta_{\text{eucl}}=0.0021$, which
would indicate an almost perfectly tree-like space; its ratio of maximum to
median token norm is $156$. Computing $\delta$ on the \emph{angular} metric
$\arccos(\cos(\cdot,\cdot))$, which is invariant to those outliers, gives
$0.2127$ --- statistically indistinguishable from BERT's $0.2114$. Flan-T5
shows the same pattern with $357\times$ norm outliers. We recommend reporting
angular $\delta$ alongside Euclidean $\delta$ together with the norm-outlier
ratio in any work motivated by measured hyperbolicity.

\paragraph{Tuning maxima are inflated.} With $T$ trials and per-seed noise
$\sigma$, the expected maximum is inflated even for an arm identical to the
baseline: $\approx 1.83$ here (STS-B/BERT, $\sigma=1.19$) and $\approx 0.83$
(STS-B/TinyLlama, $\sigma=0.54$). Several apparent leads in this project were
smaller than this and did not survive confirmation. Relatedly, the noise floor
must be measured for the configuration being confirmed: we initially assumed
CoLA's $\sigma=0.81$ from a layer-12 probe, where the true value at layer 8 is
$2.09$.

The RoBERTa campaign gives the sharpest illustration, because the inflation
there is large enough to \emph{reverse the ranking}. Mean drop from best-tuning
score to the 15-seed confirmation mean is $2.08$ MCC on CoLA and $0.48$
Spearman on STS-B, and no arm's rank is preserved: \texttt{density\_fix} is
first of five at the tuning maximum ($56.77$) and last of five over 15 seeds
($53.19$). Any protocol that reports the tuning maximum would have concluded
the opposite of what we report in \S\ref{sec:calib}.

\paragraph{Layer selection overfits.} Selecting a backbone layer on one task
does not transfer: layer 8 beat layer 12 by $+4.62$ MCC on CoLA but by only
$+0.25$ on STS-B, and lost on RTE ($-0.36$) and MRPC ($-0.46$).

\section{Notes for reproducers}
\label{app:defects}

Three details of the released implementation cost us time and are recorded here
so they do not cost anyone else's. None affects the conclusions of the original
paper.

\begin{itemize}
\itemsep1pt
\item \texttt{jk\_mode=lstm} is accepted by the argument parser but raises at
      runtime, and the advertised options \texttt{mean} and \texttt{none} are
      not implemented. The effective space is $\{$\texttt{cat},
      \texttt{max}$\}$.
\item Thresholding Poincar\'e \emph{distance} at an absolute radius yields an
      empty graph on real features. BERT token norms ($\approx 14.7$) saturate
      $\exp^c_0$ to radius $1.0$ exactly in \texttt{float32}, and real pairwise
      geodesic distances land in $[8.85, 11.76]$, so any absolute threshold
      below $\approx 9$ gives zero edges. The diagnostic signature is a score
      that is bit-identical across an entire threshold grid. We use a quantile
      threshold throughout for this reason.
\item The logged edge-density statistic counts self-loops in its numerator but
      excludes the diagonal from its denominator, so it reports
      $|E|/(n(n{-}1))$ with $|E|$ including the $n$ self-loops. A complete
      graph therefore logs $n/(n{-}1) > 1$ and an \emph{empty} graph logs
      $1/(n{-}1) \approx 0.10$ at CoLA's typical length rather than $0$. This
      is a reporting defect only --- no training path consumes the statistic,
      so no score is affected --- but it is a trap for exactly the audit this
      paper performs. Two consequences are worth recording. Our
      \texttt{no\_graph} control logs density $0.108$ on BERT/CoLA, which looks
      like a live graph and is in fact the empty-graph signature $1/(n{-}1)$;
      and the corrected BERT baseline density, $0.2074 - 0.10 \approx 0.104$,
      independently reproduces the $0.103$ measured directly in
      Table~\ref{tab:density}, which is what confirms the correction. Any
      density we quote is re-measured with the standalone instrument described
      in Table~\ref{tab:density}, never taken from training telemetry.
\end{itemize}

\end{document}
